\documentclass[10pt,a4paper]{amsart}
\usepackage[margin=19mm]{geometry}
\usepackage{amsmath,amssymb,mathtools}
\usepackage[scr=rsfs,cal=euler]{mathalfa}
\usepackage{xfrac}
\usepackage[unicode,hidelinks,bookmarks=false]{hyperref}
\newcommand{\ie}{i.e.\ }
\newcommand{\cf}{cf.\ }
\newcommand{\resp}{respectively\ }
\newcommand{\Subsection}{\subsection}
\providecommand{\nobreakdash}{\nobreak\hbox{-}}
\setlength{\emergencystretch}{2em}
\multlinegap0pt
\usepackage{amssymb}
\usepackage{float}
\usepackage{tikz}
\usepackage[shortlabels]{enumitem}
\newtheorem{theorem}{Theorem}
\title{On Matrices Whose Distinct Eigenvalues Are Fully Captured by Quotient Matrices}
\author{{ Bilal Ahmad Rather$^{a}$}}
\date{Source: arXiv:2604.03194v1 (CC BY 4.0)}
\begin{document}
\maketitle

\noindent\emph{Source and licence.} On Matrices Whose Distinct Eigenvalues Are Fully Captured by Quotient Matrices. Version arXiv:2604.03194v1 (CC BY 4.0), distributed by arXiv under the Creative Commons Attribution 4.0 International licence, http://creativecommons.org/licenses/by/4.0/, as recorded in the arXiv licence metadata for this exact version and shown on the abstract page https://arxiv.org/abs/2604.03194v1. Full licence notice: \texttt{licenses/arxiv-2604.03194-CC-BY-4.0.txt}.\par\smallskip

\noindent\textit{Editorial scope.} Counted unit: the article's theorem thm:n-by-n-two-eigs (source0.tex lines 718--731). The article's complete original proof of it is delivered with it (source0.tex lines 732--844, one \texttt{proof} environment of 113 lines). No mathematics is written, repaired or re-derived: the statement and the proof are the article's own lines, in the article's own notation, and no retained line is substituted or re-flowed.  No further source-local context is needed: every cross-reference of every retained line resolves inside the two delivered spans.  Nothing was omitted from any retained span, so the package carries no omission record.  No retained line contains a citation, so the wrapper carries no bibliography and no bibliography entry.  No retained line contains a figure environment, an image-inclusion command or drawing code, so no image is delivered and the package declares no asset dependency.  Editorial formatting only, in addition: an AMS \texttt{amsart} preamble that restates the article's own declarations and macros used by the retained lines, namely the environments \texttt{theorem} on separate counters, together with the standard packages those lines need.  The retained environments print in this document as Theorem 1 in order of appearance, whereas the article prints the same items with the numbering of its own full document.  The complete native source of the article as distributed in the arXiv e-print for this exact version is bundled unchanged as \texttt{sources/197791/source0.tex}.
\begin{theorem}\label{thm:n-by-n-two-eigs}
	Let $n\ge 3$ and let $Q=\begin{bmatrix} c_{11} & c_{12} \\ c_{21} & c_{22} \end{bmatrix}$
	be a $2\times 2$ matrix with two distinct eigenvalues $ \alpha\neq\beta.$ If $\alpha$ is the eigenvalue of matrix $M_{n}$ with multiplicity $n-1$. Then  with equitable  partition $\pi=\bigl\{\{1\},\{2,3,\dots,n\}\bigr\}$, $M$ is given by
	\begin{equation}\label{eq:Mn-two-eigs}
		M
		=
		\begin{bmatrix}
			c_{11} & \dfrac{c_{12}}{\,n-1\,}\,\mathbf{1}_{1\times (n-1)}\\
			c_{21}\,\mathbf{1}_{(n-1)\times 1} & \alpha I_{n-1} + \dfrac{c_{22}-\alpha}{\,n-1\,}J_{n-1}
		\end{bmatrix},
	\end{equation}
	where $\mathbf{1}_{r\times s}$ is the $r\times s$ all--ones matrix, $I_{n-1}$ is the identity, and
	$J_{n-1}$ is the $(n-1)\times(n-1)$ all--ones matrix.
\end{theorem}

\begin{proof}
	Write $M\in \mathbb{C}^{n\times n}$ in block form with respect to the partition $\pi=\{\{1\},\{2,\dots,n\}\}$ as
	$$
	M
	=
	\begin{bmatrix}
		m_{11} & r^\mathsf{T}\\
		s & B
	\end{bmatrix},
	$$
	where $r,s\in\mathbb{C}^{n-1}$ and $B\in\mathbb{C}^{(n-1)\times(n-1)}$. For $\pi$ to be equitable with quotient $Q=\bigl(c_{ij}\bigr)_{1\le i,j\le 2}$, the row-sums from each cell to each cell must be constant.  From cell $C_1=\{1\}$ to $C_1$, the sum is $m_{11}$, so we must have $m_{11} = c_{11}.$  From $C_1$ to $C_2$,  the sum is  $\sum_{j=2}^n m_{1j} = \sum_{j=1}^{n-1} r_j = c_{12}.$ From $C_2$ to $C_1$, for every $i\in\{2,\dots,n\}$, the entry $m_{i1}$ must be the same, say $c_{21}$, and we get  $m_{i1}=c_{21}$ for all $i\ge 2,$ that is, $s=c_{21}\mathbf{1}$.  From $C_2$ to $C_2$, for each $i\in\{2,\dots,n\}$, the row sum across columns $2,\dots,n$ must be constant,
		$$
		\sum_{j=2}^n B_{ij} = c_{22}\quad\text{for all }i\ge 2.
		$$	
	The choices in \eqref{eq:Mn-two-eigs} now leave us with
	$$
	m_{11}=c_{11},\qquad
	r=\frac{c_{12}}{\,n-1\,}\mathbf{1}_{n-1},\quad
	B = \alpha I_{n-1} + \frac{c_{22}-\alpha}{\,n-1\,}J_{n-1}.
	$$
	So, we have
	$$
	\sum_{j=2}^n m_{1j} 
	= \sum_{j=1}^{n-1}\frac{c_{12}}{\,n-1\,}
	= c_{12}.
	$$
	Further, for each $i\ge 2$, 
	$$
	m_{i1}=c_{21},\quad
	\sum_{j=2}^n B_{ij}
	= \alpha + \frac{c_{22}-\alpha}{\,n-1\,}\cdot (n-1) 
	= \alpha + (c_{22}-\alpha)
	= c_{22}.
	$$
	As a result, the row sum from one cell to the next are $\begin{bmatrix}
		c_{11} & c_{12}\\
		c_{21} & c_{22}
	\end{bmatrix}
	=Q$. In this case, $\pi$ is equitable, and its equitable quotient matrix is $Q$.
	
	 	Let $\mathbf{1}=\mathbf{1}_{n-1}=\begin{bmatrix}1\\\vdots\\1\end{bmatrix}\in\mathbb{C}^{n-1},$ and $	u =\begin{bmatrix}0\\ \mathbf{1}\end{bmatrix}\in\mathbb{C}^{n}.$ Then  the partition subspace spanned by $e_1$ and $u$ is
	$$
	\mathcal{S}_\pi=\{(a,b,\dots,b)^\mathsf{T}:a,b\in\mathbb{C}\}
	$$
	 Its orthogonal complement of dimension $n-2$ with respect to the standard inner product is
	$$
	\mathcal{S}_\pi^\perp
	=
	\{(0,x_2,\dots,x_n)^\mathsf{T} : x_2+\dots+x_n=0\}.
	$$
	
	\smallskip
	By standard theory of equitable partitions, the restriction of $M$ to $\mathcal{S}_\pi$ in the basis $\{e_1,u\}$ is represented by the quotient matrix $Q$. Thus, the eigenvalues of $M|_{\mathcal{S}_\pi}$ are exactly $\alpha$ and $\beta$, which are same as the eigenvalues of $Q$.
	
	Let's look at what happens on $\mathcal{S}_\pi^\perp$.Let $x=(0,y)^\mathsf{T}$ where $y\in\mathbb{C}^{n-1}$ and $\mathbf{1}^\mathsf{T}y=0$. After that, we have
	$$
	Mx
	=
	\begin{bmatrix}
		c_{11} & r^\mathsf{T}\\
		s & B
	\end{bmatrix}
	\begin{bmatrix}0\\y\end{bmatrix}
	=
	\begin{bmatrix}
		r^\mathsf{T}y\\
		By
	\end{bmatrix}.
	$$
	There is no way for $r^\mathsf{T}y=0$ since $r$ is a scalar multiple of $\mathbf{1}_{n-1}$ and $y$ has no sum. Now,
	$$
	By
	=
	\left(\alpha I_{n-1} + \frac{c_{22}-\alpha}{\,n-1\,}J_{n-1}\right)y
	=
	\alpha y + \frac{c_{22}-\alpha}{\,n-1\,} J_{n-1}y.
	$$
	However, $J_{n-1}y = (\mathbf{1}^\mathsf{T}y)\mathbf{1}=0$, which means that $By=\alpha y$ and we get
	$$
	Mx = \begin{bmatrix}0\\ \alpha y\end{bmatrix} = \alpha x.
	$$
	So, each vector in $\mathcal{S}_\pi^\perp$ has an eigenvalue of $\alpha$ and is an eigenvector of $M$. Since $\dim\mathcal{S}_\pi^\perp = n-2$, this contributes $\alpha$ with geometric (and hence algebraic) multiplicity at least $n-2$. Together with the eigenvalue $\alpha$ coming from the restriction $M|_{\mathcal{S}_\pi}$ (because $Q$ has eigenvalues $\alpha,\beta$), we see that $\alpha$  has algebraic multiplicity $n-1,$	and $\beta$ appears once, coming from the other eigenvalue of $Q$. Hence, the spectrum of $M$ is  $\operatorname{spec}(M)=\{\beta,\alpha,\dots,\alpha\}.$
	
	\medskip
	Next, we shoe that $M$ must be necessarily of the form \eqref{eq:Mn-two-eigs}. Suppose $M\in\mathbb{C}^{n\times n}$ has exactly two distinct eigenvalues $\alpha$ and $\beta$ with algebraic multiplicities $n-1$ and $1$, and that the partition $\pi=\{\{1\},\{2,\dots,n\}\}$ is equitable with quotient $Q$. Let $S_\pi$ and $\mathcal{S}_\pi^\perp$ be as above, and assume $\mathcal{S}_\pi^\perp \subseteq E_\alpha.$  Recall,   $M$ as before:
	$$
	M=\begin{bmatrix}m_{11} & r^\mathsf{T}\\ s & B\end{bmatrix}.
	$$
	Equitability condition of $\pi$  forces us
	$$
	m_{11}=c_{11},\sum_{j=1}^{n-1}r_j=c_{12},
	s=c_{21}\mathbf{1},\quad \sum_{j=2}^n B_{ij}=c_{22}\ \text{for all } i.
	$$
	Now, take $x=(0,y)^\mathsf{T}\in\mathcal{S}_\pi^\perp$, so $\mathbf{1}^\mathsf{T}y=0$ and $Mx=\alpha x$.
	As before
	$$
	Mx=\begin{bmatrix}r^\mathsf{T}y\\ By\end{bmatrix}=\alpha\begin{bmatrix}0\\y\end{bmatrix}.
	$$
	So, $r^\mathsf{T}y=0$, for all $y$ with $\mathbf{1}^\mathsf{T}y=0$. This implies $r$ lies in the span of $\mathbf{1}$, that is, $r=\delta\mathbf{1}$ for some scalar $\delta$, and from the row-sum condition we obtain
	$$
	\delta = \frac{c_{12}}{\,n-1\,}.
	$$
	Similarly, $By=\alpha y$, for all $y$ with $\mathbf{1}^\mathsf{T}y=0$. This means that
	$$
	(B-\alpha I_{n-1})y=0\quad\text{for all }y\perp \mathbf{1}.
	$$
	So, it follows that $B-\alpha I_{n-1}$ has image contained in $\operatorname{span}\{\mathbf{1}\}$ and kernel containing $\{\mathbf{1}\}^\perp$. Thus $B-\alpha I_{n-1}$ is a rank--$1$ matrix of the form
	$B-\alpha I_{n-1} = u\mathbf{1}^\mathsf{T}$,	for some $u\in\mathbb{C}^{n-1}$. The equitability condition on row-sums $\sum_{j}B_{ij}=c_{22}$ gives us
	$$
	\alpha + (n-1)u_i = c_{22}\quad\text{for all }i.
	$$
	So, $u_i$ is constant, say $u_i=\gamma$ for all $i$, and $\gamma = \frac{c_{22}-\alpha}{\,n-1\,}.$	Thus, $B = \alpha I_{n-1} + \dfrac{c_{22}-\alpha}{\,n-1\,}J_{n-1}$, and we recover \eqref{eq:Mn-two-eigs}. This shows $M$ must have the claimed form.  That completes the proof.
\end{proof}

\end{document}
